% GENERATED FILE. Edit canonical lessons, then run npm run generate:books.
% canonical-source-hash: fnv1a64:efad6dbd01265c16
% canonical-lessons: BN-R01-tumi-again, BN-R01-ta-u-again, BN-R01-ya-again, BN-R01-chandrabindu-again, BN-R01-hyan-again, BN-R01-ami-again, BN-R01-kemon-again, BN-R01-bhalo-again, BN-R01-practice

\chapter{The Distant Band}
\label{ch:bn-distant-band}
\hlchaptermodality{26}

\begin{chapteropening}
\textbf{By the end of this chapter:} \emph{I can retrieve, with nothing uncovered, the words and letters this book taught between ninety-eight and a hundred and four lessons ago — and I can tell the difference between a word I have lost and a word I now know to go back for.}

Four chapter-closing practice lessons, untouched for the length of a whole book, are answered in one sitting; nothing in the chapter is new, and that is the measurement.
\end{chapteropening}

\section[Practice]{\bn{তুমি}, at the far end}

\label{lesson:BN-R01-tumi-again}



\begin{quote}
Before anything is uncovered: say \emph{you} to a friend, \emph{you} to a stranger, and \emph{what}.
\end{quote}

\subsection*{Your turn — the distant band}
This chapter has one job, and it is not to teach you anything.

Everything asked for here was first taught between a hundred and a hundred and four lessons ago. That distance is deliberate. A thing recalled after a day is being remembered; a thing recalled after a hundred lessons is being \textbf{owned}, and the only way to find out which you have is to reach that far and see what comes back.

Nothing new arrives in this chapter. If a word does not come, do not treat that as a failure — treat it as the measurement working.

Say these aloud:

\begin{itemize}
\raggedright
  \item \emph{you}, to a friend — \emph{tumi}
  \item \emph{you}, to someone you have only now met — \emph{āpni}
  \item \emph{what} — \emph{ki}
  \item the whole question, \textquotedblleft{}what's your name?\textquotedblright{} — \emph{tomār nām ki?}
  \item the same question, respectfully — \emph{āpnār nām ki?}
\end{itemize}

Three of those five you have not been asked for in eighty lessons.

\subsection*{Writing — from sound}
Cover every Bengali model in this book. Write \emph{tumi} from the sound alone, then write the vowel \textbf{\bn{ই}} standing on its own.

Write these out:

\begin{itemize}
\raggedright
  \item \bn{তুমি}
  \item \bn{ই}
\end{itemize}

\begin{tcolorbox}[breakable,colback=teal!4,colframe=teal!35!black,title={Before you move on}]
Which sign in \textbf{\bn{তুমি}} hangs below the line? (\textbf{\bn{ু}}.) And which \textquotedblleft{}you\textquotedblright{} would you use with a shopkeeper you do not know? (\textbf{\bn{আপনি}}.)
\end{tcolorbox}
\section[Practice]{\bn{নাম}, and two pieces from the signs chapter}

\label{lesson:BN-R01-ta-u-again}



\begin{quote}
Say \emph{name}, then \emph{my}, then put them together.
\end{quote}

\subsection*{Your turn — the distant band}
Say these aloud:

\begin{itemize}
\raggedright
  \item \emph{name} — \emph{nām}
  \item \emph{my} — \emph{āmār}
  \item \textquotedblleft{}my name is —\textquotedblright{} with your own name in it
  \item what Bengali does \emph{not} put in that sentence (\textbf{no \textquotedblleft{}is\textquotedblright{}})
\end{itemize}

The missing verb is the part that decays first, because English keeps insisting it belongs there. Say the sentence once more without it.

\subsection*{Writing — from sound}
Cover every model. Write the dental consonant, then the sign that hangs beneath a letter, then the word for \emph{name}.

Write these out:

\begin{itemize}
\raggedright
  \item \bn{ত}
  \item \bn{ু}
  \item \bn{নাম}
\end{itemize}

\begin{tcolorbox}[breakable,colback=teal!4,colframe=teal!35!black,title={Before you move on}]
Is Bengali's \textbf{\bn{ত}} made at the teeth or behind them? (\textbf{At the teeth}.) And how many pieces are in \textbf{\bn{নাম}}? (\textbf{Three}.)
\end{tcolorbox}
\section[Practice]{\bn{য}, and the thanks it hides in}

\label{lesson:BN-R01-ya-again}



\begin{quote}
Say \emph{thank you}, and then the thing a Bengali speaker is more likely to say instead of it.
\end{quote}

\subsection*{Your turn — the distant band}
Say these aloud:

\begin{itemize}
\raggedright
  \item \emph{thank you} — \emph{dhônyobad}
  \item \textquotedblleft{}it's no matter / you're welcome\textquotedblright{} — \emph{kono bæpār nā}
  \item what \textbf{\bn{য}} sounds like standing alone (\emph{\textbf{j}})
  \item what it sounds like hung under another letter (\textbf{nothing} — it bends the vowel to \textbf{æ})
\end{itemize}

Both of the words above carry that hung \textbf{\bn{য}}, and it is the same job in each: the \textbf{ny} in \emph{dhônyobad} and the \textbf{byæ} in \emph{bæpār}. Two words from two different chapters, one piece of machinery.

\subsection*{Writing — from sound}
Cover every model. Write the letter that is written \emph{y} and said \emph{j}, then the mark that kills an inherent vowel, then hang the first under a \textbf{\bn{হ}}.

Write these out:

\begin{itemize}
\raggedright
  \item \bn{য}
  \item \bn{্}
  \item \bn{হ্য}
\end{itemize}

\begin{tcolorbox}[breakable,colback=teal!4,colframe=teal!35!black,title={Before you move on}]
What is the scribes' name for a \textbf{\bn{য}} hung below? (\emph{\textbf{jô-phôla}}.) And what does \textbf{\bn{্}} do? (\textbf{Deletes the inherent vowel}.)
\end{tcolorbox}
\section[Practice]{\bn{ঁ}, from above the line}

\label{lesson:BN-R01-chandrabindu-again}



\begin{quote}
Say \emph{yes} and \emph{no}.
\end{quote}

\subsection*{Your turn — the distant band}
Say these aloud:

\begin{itemize}
\raggedright
  \item \emph{yes} — \emph{hyã}, through the nose
  \item \emph{no} — \emph{nā}, not through the nose
  \item the number five — \emph{pā̃ch}, and notice the same nasal
  \item where on the line the mark that writes it sits (\textbf{above})
\end{itemize}

Three words in this whole book carry that crescent, and they are spread across three chapters. That is exactly the shape of thing a reader loses: rare enough to go, common enough to matter.

Say \emph{hyã} and \emph{nā} one after another until the difference is in the nose and not in the effort.

\subsection*{Writing — from sound}
Cover every model. Write the crescent-with-a-dot, then put it on \textbf{\bn{না}}.

Write these out:

\begin{itemize}
\raggedright
  \item \bn{ঁ}
  \item \bn{নাঁ}
\end{itemize}

\begin{tcolorbox}[breakable,colback=teal!4,colframe=teal!35!black,title={Before you move on}]
Is \textbf{\bn{ঁ}} a consonant? (\textbf{No}.) What does it change? (\textbf{The vowel} — it goes through the nose.) And is it the only mark above the head-line? (\textbf{Yes}.)
\end{tcolorbox}
\section[Practice]{\bn{হ্যাঁ} and \bn{নমস্কার}, side by side}

\label{lesson:BN-R01-hyan-again}



\begin{quote}
Say the greeting, with the hands.
\end{quote}

\subsection*{Your turn — the distant band}
Say these aloud:

\begin{itemize}
\raggedright
  \item hello, and goodbye, with one word — \emph{nômoshkar}
  \item \emph{yes} — \emph{hyã}
  \item how many pieces are in each (\textbf{seven}, and \textbf{five})
\end{itemize}

These two words are the bookends of the reading strand. The greeting was the first thing you could decode and it took nine lessons of letters to earn. \emph{Yes} came three chapters later and cost three.

The difference is not that the second word is simpler. It is that by the time it arrived you already owned four of its five pieces — which is the argument for buying letters early and spending them often, made in two words.

\subsection*{Writing — from sound}
Cover every model. Write \emph{yes}, then the breath-letter alone, then the last consonant of the greeting.

Write these out:

\begin{itemize}
\raggedright
  \item \bn{হ্যাঁ}
  \item \bn{হ}
  \item \bn{র}
\end{itemize}

\begin{tcolorbox}[breakable,colback=teal!4,colframe=teal!35!black,title={Before you move on}]
Which piece of \textbf{\bn{হ্যাঁ}} makes no sound? (\textbf{\bn{য}}.) And which word took nine lessons to become readable? (\textbf{\bn{নমস্কার}}.)
\end{tcolorbox}
\section[Practice]{\bn{আমি}, and the sentence that can drop it}

\label{lesson:BN-R01-ami-again}



\begin{quote}
Say \emph{I}.
\end{quote}

\subsection*{Your turn — the distant band}
Say these aloud:

\begin{itemize}
\raggedright
  \item \emph{I} — \emph{āmi}
  \item \textquotedblleft{}my name is —\textquotedblright{}
  \item \textquotedblleft{}pleased to meet you\textquotedblright{} — \emph{ālāp kore bhālo lāglo}
  \item \textquotedblleft{}I am well\textquotedblright{} \textbf{without} the pronoun (\emph{bhālo āchhi})
\end{itemize}

That last one is the test. The pronoun is carried by the \textbf{\bn{ি}} on the end of the verb, so a speaker who drops \textbf{\bn{আমি}} loses nothing at all — and a learner who cannot drop it is still translating rather than speaking.

Say the polite formula once more, whole, at speed. It is the longest thing this book asks anyone to say from memory.

\subsection*{Writing — from sound}
Cover every model. Write \emph{I} from the sound alone.

\begin{itemize}
\raggedright
  \item \emph{Write it:} \bn{আমি}
\end{itemize}

\begin{tcolorbox}[breakable,colback=teal!4,colframe=teal!35!black,title={Before you move on}]
Which piece of the verb carries \textquotedblleft{}I\textquotedblright{}? (\textbf{The \bn{ি} ending}.) So is \textbf{\bn{আমি}} required? (\textbf{No} — it is emphasis.)
\end{tcolorbox}
\section[Practice]{\bn{কেমন}, and the whole exchange}

\label{lesson:BN-R01-kemon-again}



\begin{quote}
Say \emph{how}.
\end{quote}

\subsection*{Your turn — the distant band}
Say these aloud:

\begin{itemize}
\raggedright
  \item \emph{how} — \emph{kēmon}
  \item \textquotedblleft{}how are you?\textquotedblright{} to a friend — \emph{tumi kēmon āchho?}
  \item the same to a stranger — \emph{āpni kēmon āchhen?}
  \item \textquotedblleft{}okay / I see / go on\textquotedblright{} — \emph{āchchhā}
\end{itemize}

Two of those four change one ending and nothing else. If the respectful form did not come, it is because the familiar one is the one you have practised, and a question you can only ask a friend is half a question.

Run the exchange with yourself, both sides, twice.

\subsection*{Writing — from sound}
Cover every model. Write \emph{how} from the sound alone.

\begin{itemize}
\raggedright
  \item \emph{Write it:} \bn{কেমন}
\end{itemize}

\begin{tcolorbox}[breakable,colback=teal!4,colframe=teal!35!black,title={Before you move on}]
What changes between \emph{āchho} and \emph{āchhen}? (\textbf{Only the ending} — who you are speaking to.) And which sign in \textbf{\bn{কেমন}} goes in front of its consonant? (\textbf{\bn{ে}}.)
\end{tcolorbox}
\section[Practice]{\bn{ভালো}, and the reply it lives in}

\label{lesson:BN-R01-bhalo-again}



\begin{quote}
Say \emph{good}.
\end{quote}

\subsection*{Your turn — the distant band}
Say these aloud:

\begin{itemize}
\raggedright
  \item \emph{good, well} — \emph{bhālo}
  \item \textquotedblleft{}I am well\textquotedblright{} — \emph{āmi bhālo āchhi}
  \item the goodbye that literally means \textquotedblleft{}I'll come again\textquotedblright{} — \emph{āshi}
  \item why a Bengali speaker prefers that goodbye (\textbf{it promises a return})
\end{itemize}

\textbf{\bn{ভালো}} has appeared in more lessons than any other word in this book, which makes it the fair test in the other direction: not \textquotedblleft{}did a rare thing survive?\textquotedblright{} but \textquotedblleft{}did constant exposure actually put a thing where you can get at it?\textquotedblright{}

If it came back instantly, that is a result worth having. Frequency is supposed to do that, and it is worth knowing when it worked.

\subsection*{Writing — from sound}
Cover every model. Write \emph{good}, then \emph{am}, then the two together with \emph{I} in front.

Write these out:

\begin{itemize}
\raggedright
  \item \bn{ভালো}
  \item \bn{আছি}
  \item \bn{আমি} \bn{ভালো} \bn{আছি}
\end{itemize}

\begin{tcolorbox}[breakable,colback=teal!4,colframe=teal!35!black,title={Before you move on}]
Which sign wraps around the \textbf{\bn{ল}} in \textbf{\bn{ভালো}}? (\textbf{\bn{ো}}, on both sides.) And what does \emph{āshi} literally promise? (\textbf{That you will come again}.)
\end{tcolorbox}
\section[Practice]{The band closed}

\label{lesson:BN-R01-practice}



\begin{quote}
Greet someone, introduce yourself, and say goodbye — all before anything is uncovered.
\end{quote}

\subsection*{Your turn — the near edge of the band}
Four chapters of this book ended in a practice lesson, and each of those lessons has been sitting untouched for the length of a whole book. Answer all four now, in order, out loud:

Say these aloud:

\begin{itemize}
\raggedright
  \item the five greeting words — \emph{nômoshkar}, \emph{dhônyobad}, \emph{hyã / nā}, \emph{āchchhā}, \emph{āshi}
  \item the introduction exchange, both halves
  \item the how-are-you exchange, both halves
  \item the two farewells — \emph{ābār dækhā hôbe}, \emph{kāl dækhā hôbe}
\end{itemize}

That is the whole speaking spine of the book in four breaths.

Now the thing worth saying plainly about this chapter. Nothing in it was new. Every item was first taught between ninety-eight and a hundred and four lessons ago, and the reason to gather them here rather than to press on is that \textbf{a curriculum that only ever adds is not teaching, it is delivering.} What you could still reach for today is what you actually have.

Anything that did not come back is not lost. It is located.

\subsection*{Writing — from sound}
Cover every model in the book. Write \emph{I}, then the consonant that was read beside it.

Write these out:

\begin{itemize}
\raggedright
  \item \bn{আমি}
  \item \bn{ল}
\end{itemize}

\begin{tcolorbox}[breakable,colback=teal!4,colframe=teal!35!black,title={Before you move on}]
How many new things did this chapter teach? (\textbf{None}.) How far back did it reach? (\textbf{About a hundred lessons}.) And is a word you could not retrieve today a word you have lost? (\textbf{No} — it is a word you now know to go back for.)
\end{tcolorbox}
